% ============================================
% Supplementary Note 7: Interpretation of Pairwise Selection Correlation
% ============================================
\subsection*{Supplementary Note 7: Interpretation of pairwise selection correlation}
\label{sec:suppl:note7}

\paragraph{\rev{Metric interpretation.}}
\rev{The pairwise selection-correlation metric is defined in Supplementary Methods. Positive $\bar{\rho}_{\text{same}}$ indicates that species with the same parental-affiliation label share fates more often than they have different fates, whereas negative $\bar{\rho}_{\text{cross}}$ indicates that cross-affiliation species pairs have different fates more often than shared fates. Thus, $\Delta=\bar{\rho}_{\text{same}}-\bar{\rho}_{\text{cross}}$ measures how strongly survival or extinction becomes correlated within assembled parental communities relative to pairs drawn from different parental communities.}

\paragraph{\rev{Simulation evidence.}}
\rev{In simulations (Extended Data Fig.~6a--d), same-parent pairs remained positively correlated, whereas cross-community pairs became increasingly negative as the interaction-strength parameter increased. At low interaction-strength parameter values ($\mu = 0.3$), both within-community and cross-community pairs show positive correlations with only a small difference ($\Delta = 0.05$; two-sided permutation test, $p < 0.001$; Extended Data Fig.~6a). As the interaction-strength parameter increases, the difference becomes pronounced. At $\mu = 0.6$, within-community pairs show positive correlation while cross-community pairs become negative ($\Delta = 0.59$, $p < 0.001$; Extended Data Fig.~6b). At $\mu = 0.8$, the separation is larger ($\Delta = 0.94$, $p < 0.001$; Extended Data Fig.~6c).}

\paragraph{\rev{Experimental evidence.}}
\rev{In experiments (Extended Data Fig.~7), the same-versus-cross separation was statistically significant in Base ($n=83$, $\Delta=0.235$, $p < 0.001$) and Nutr$+$ ($n=83$, $\Delta=0.141$, $p < 0.001$), but not in Nutr$-$ ($n=90$, $\Delta=0.016$, $p=0.307$). The separation was largest in Base, and the significant Nutr$+$ separation indicates parental-affiliation-correlated species fates under a condition with higher failed-invasion frequency and stronger environmental feedbacks.}

\paragraph{\rev{Biological interpretation.}}
\rev{Together, these patterns indicate that assembly history structures species-fate concordance during coalescence: species from the same assembled parental community tend to persist or be lost together, whereas species from different parental communities become increasingly anti-concordant as interactions strengthen. This provides species-fate-level support for interpreting Dominance outcomes as community-level selection~\citep{Tikhonov2016,DiazColunga2022,Mansour2018,Rillig2015}.}

\paragraph{\rev{Relationship to invasion fitness.}}
\rev{This metric can also be interpreted through the invasion-fitness framework. Related pairwise invasion-resistance assays are summarized in Supplementary Note~4. In a gLV model, invasion fitness is the per-capita growth rate of a rare species introduced into a resident community background. Positive invasion fitness predicts establishment, whereas negative invasion fitness predicts exclusion~\citep{Hu2025,Kurkjian2021}. The 95:5 pairwise invasion assays provide an empirical two-species approximation of this idea, while the pairwise fate-concordance metric asks whether analogous persistence outcomes become correlated across many species from the same assembled parental community during coalescence. Consistent with this interpretation, an auxiliary gLV analysis showed that excess same-parent shared fate, measured relative to a parental-affiliation-shuffled null, increased with the interaction-strength parameter $\mu$ across the full simulated sweep ($\mu = 0.05$--$1.20$; Pearson $r = 0.870$, $p = 3.2 \times 10^{-8}$; Supplementary Fig.~36). This auxiliary within-parent metric measures the same parental-affiliation-correlated shared-fate signal as $\Delta=\bar{\rho}_{\text{same}}-\bar{\rho}_{\text{cross}}$, but compares same-parent pairs to a parental-affiliation-shuffled null rather than explicitly subtracting cross-community pairs.}
